\documentclass[11pt,a4paper]{article}
\usepackage[T1]{fontenc}
\usepackage[utf8]{inputenc}
\usepackage{amsmath,amssymb,amsthm}
\usepackage[margin=1in]{geometry}
\usepackage[colorlinks=true,linkcolor=blue,urlcolor=blue]{hyperref}
\theoremstyle{plain}
\newtheorem{theorem}{Theorem}
\newtheorem{lemma}[theorem]{Lemma}
\newtheorem{proposition}[theorem]{Proposition}
\newtheorem{corollary}[theorem]{Corollary}
\theoremstyle{definition}
\newtheorem{remark}[theorem]{Remark}

\title{No column above the one before it in every row: OEIS A212929}
\author{Adrian Perez Fontelles\\ \small Independent researcher}
\date{7 September 2026}
\begin{document}
\maketitle

\begin{abstract}
OEIS A212929 carries an empirical recurrence of order $15$ contributed by R.~H.~Hardin. It is
true, and it is decidable rather than empirical. The entry's condition quantifies over all rows at once, so it is not decided by any
bounded window; but it is a conjunction of independent per-row facts, and one bit per
adjacent column pair carries it; the
admissible windows are the vertices of a finite digraph and the arrays counted are exactly the
walks in it. Hence $a(n)$ is a walk count on $S=15$ vertices and is $C$-finite, and whether the
conjectured recurrence annihilates it is settled by a finite exact computation, with the
Cayley--Hamilton theorem bounding the work.
\end{abstract}

\noindent\small 2020 Mathematics Subject Classification. 05A15, 05B45, 15B36.\normalsize

\section{The conjecture}

OEIS A212929 is ``Number of n X 7 0..6 arrays with no column j greater than column j-1 in all rows.''. It has offset $1$ and begins
\[
1716, 163050236504, 330676305397366860, 372271852759802667094256, 346763870885473008105706659876, 300496032699305814168049256913886664, 252848108054080530934326716965286554552380,\ \dots
\]
The entry states, as a conjecture contributed by its author and never marked settled:
\begin{quote}
Empirical: a(n) = 1556535*a(n-1) -792039717385*a(n-2) +178076842413452815*a(n-3) -20358575698888213461493*a(n-4) +1266623824167234235679318755*a(n-5) -44047909952609917776536882759565*a(n-6) +864953744390957718129128593210421035*a(n-7) -9623052428363321062071131847418214552931*a(n-8) +59322745910799635132647204742065665730348485*a(n-9) -190454450982198883185256207135935908957968246875*a(n-10) +287023164786865836968027438435034735344099822518125*a(n-11) -193456826915294502754189697430417253206202705301574375*a(n-12) +48193709324773967360630681370925021204404072468778790625*a(n-13) -1968890430914107376575740256318014532178839614309088359375*a(n-14) +1920889891583479054710502398832778506621251127702072265625*a(n-15)
\end{quote}
As of the ``Last modified'' line on the live entry (Oct 03 2025 16:40:26, revision 12) this is still
recorded as empirical, and nothing on the entry records it as proved.

\section{The count is a walk count}

The entry counts arrays of $7$ columns whose rows are arbitrary over the alphabet $\{0,1,\dots,6\}$, and asks that no column $j$ be greater than column $j-1$ \emph{in all rows}.

That is not a condition on any bounded window: ``in all rows'' quantifies over the whole array, and a row arriving at the very end can still decide it. What makes it a walk anyway is that it is a conjunction of independent facts about single rows. For each of the $6$ adjacent column pairs keep one bit, recording whether every row so far has had column $j$ greater than column $j-1$. A new row can only clear bits, never set one, and the array is admissible exactly when every bit has been cleared by the time the array ends.

So the state is that bit mask and nothing else --- at most $2^{6}$ of them, however large the alphabet is, which is why entries whose terms run to twenty digits have models this small. Rows with the same bit pattern act identically, so each row contributes one edge and the multiplicity does the counting. Merging vertices with identical futures leaves $S=15$, and with $M$ the adjacency matrix and $\iota,\tau$ the vectors recording which masks may begin and end an array,
\[
a(n)\;=\;\iota^{\!\top}M^{\,n-1}\tau ,
\]
the exponent fixed by the entry's own shape and checked against its published terms. In
particular $a$ is $C$-finite.

\section{The proof}

\begin{lemma}\label{lem:crit}
Let $q(t)=t^{15}-\sum_i c_it^{15-i}$ be the characteristic polynomial of the
conjectured recurrence. The residual $a(n)-\sum_ic_ia(n-i)$ equals
$\iota^{\!\top}M^{\,j}q(M)\tau$ for the corresponding walk index $j$, so the recurrence
holds from some point on if and only if $\iota^{\!\top}M^{j}q(M)\tau=0$ for all large $j$,
and it is enough to examine $j<S$.
\end{lemma}

\begin{proof}
Factor $M^{j}$ out of $M^{j+15}-\sum_ic_iM^{j+15-i}$. For the bound, $M^{S}$
is an integer combination of $I,M,\dots,M^{S-1}$ by Cayley--Hamilton, so the numbers
$u_j=\iota^{\!\top}M^{j}q(M)\tau$ obey the monic recurrence given by the characteristic
polynomial of $M$; $S$ consecutive zeros force every later one, and the last nonzero $u_j$
pins the threshold exactly. A constant factor in front of $a$ divides out of the residual and
changes nothing here.
\end{proof}

\begin{theorem}
\[
a(n)\;=\;1556535\,a(n-1) - 792039717385\,a(n-2) + 178076842413452815\,a(n-3) - 20358575698888213461493\,a(n-4) + 1266623824167234235679318755\,a(n-5) - 44047909952609917776536882759565\,a(n-6) + 864953744390957718129128593210421035\,a(n-7) - 9623052428363321062071131847418214552931\,a(n-8) + 59322745910799635132647204742065665730348485\,a(n-9) - 190454450982198883185256207135935908957968246875\,a(n-10) + 287023164786865836968027438435034735344099822518125\,a(n-11) - 193456826915294502754189697430417253206202705301574375\,a(n-12) + 48193709324773967360630681370925021204404072468778790625\,a(n-13) - 1968890430914107376575740256318014532178839614309088359375\,a(n-14) + 1920889891583479054710502398832778506621251127702072265625\,a(n-15)
\]
for every $n>15$.
\end{theorem}

\begin{proof}
The vector $w=q(M)\tau$ was formed by $15$ matrix--vector products in exact integer
arithmetic and $\iota^{\!\top}M^{j}w$ evaluated for $j=0,1,\dots$, again exactly; those
integers vanish from the index corresponding to $n=15+1$ onwards and the run continued
until the working vector was identically zero, which settles every larger $j$ at once.
\end{proof}

\section{Verification}

Four checks, each independent of the algebra above.

First, the model reproduces all $7$ terms the entry publishes, exactly, in integer
arithmetic. This is what ties the model to the entry: the digraph is built by reading the
entry's English, and a misreading gives different counts.

Second, the reading was pinned before the digraph was written, by a program that enumerates
every array of the given shape one by one and tests, for each adjacent column pair,
whether every row really does put one above the other. That program shares no code with the transfer matrix, so an error in one does not
hide an error in the other.

Third, the conjectured recurrence was evaluated directly on the published terms, with no
matrices involved, and holds wherever the proved range and the data overlap.

Fourth, the annihilation test of Lemma~\ref{lem:crit} was carried out in exact integer
arithmetic on the whole vector rather than on a sampled prefix.

\begin{thebibliography}{9}
\bibitem{oeis} The OEIS Foundation, \emph{The On-Line Encyclopedia of Integer
Sequences}, \url{https://oeis.org}, 2026. Sequence A212929.
\bibitem{stanley} R.~P.~Stanley, \emph{Enumerative Combinatorics}, Volume 1, 2nd ed.,
Cambridge University Press, 2011. (Transfer-matrix method, Section 4.7.)
\bibitem{fs} P.~Flajolet and R.~Sedgewick, \emph{Analytic Combinatorics}, Cambridge
University Press, 2009.
\bibitem{horn} R.~A.~Horn and C.~R.~Johnson, \emph{Matrix Analysis}, 2nd ed., Cambridge
University Press, 2013.
\end{thebibliography}

\end{document}
